\scanpage{44}
\noindent\hfill atkin \mspage{4}

\noindent These are all due to LEHNER(1949a,1949b,1950) except the cases \uline{$e$}$\handwritten{>}3$ due to ATKIN(1967). They represent what are perhaps the simplest and most natural ``Ramanujan-type'' congruences. The general possibilities are not confined to the coefficients \uline{$c(n)$} and \uline{$p(n)$},and many other modular forms have Fourier coefficients* with similar p\rlap{\handwritten{\tikz[baseline=0pt,x=1em,y=1ex]{\draw (0,-.35)--(.56,-.35)--(.56,1.7)--(1.06,1.7);}}}orperties. The work

\noindent\handwritten{\rule{\linewidth}{0.4pt}}

\noindent *That is,the coefficients of the expansion of a form $F(\handwritten{\unsure{\cdot|}\tau})$ in powers of $x=\handwritten{e^{2\pi i\tau}}$.

\noindent\handwritten{\rule{\linewidth}{0.4pt}}

\noindent described in this paper was in fact conducted simultaneously with \uline{$c(n)$} and \uline{$p(n)$},the former being the simplest case and the latter an example of the most difficult type of case; but we confine ourselves to \uline{$c(n)$} from now on in the present account.

\noindent 3. \uline{Congruence}s \uline{for} $q=13$. It might have been hoped that results similar to \handwritten{(2.5)} would exist for primes $q\handwritten{>}11$,but this is(by now)extremely improbable(there are strong theoretical grounds for this assertion which is not merely based on the fact that no properties have been found).

\noindent The first results for $q\handwritten{>}11$,of an entirely different kind from those described in \handwritten{\S}2,are due to NEWMAN(1958),who proved that
\[
\handwritten{c(13^{2}n)\equiv8\,c(13n)\quad(\mathrm{mod}\;13),}
\tag*{\handwritten{(3\unsure{.|,}1)}}
\]
\[
\begin{aligned}
&\handwritten{\text{and}\quad\text{if}\quad t(n)=-c(13n)\quad\text{and}\quad p\ne13\quad\text{is prime then}}\\
&\qquad\handwritten{t(np)-t(n)t(p)+p^{11}\,t(n/p)\equiv0\,(\mathrm{mod}\;13).}
\end{aligned}
\tag*{\handwritten{$(3,2)^{*}$}}
\]

\noindent\handwritten{\rule{\linewidth}{0.4pt}}

\noindent *For any number-theoretic coefficient \uline{$t(n)$} we write $t(n)=0$ if \uline{$n$} is not an integer.

\noindent\handwritten{\rule{\linewidth}{0.4pt}}

From these Newman deduces,using the computed ``accident'' that $c(7\mathbin{.}13)\handwritten{\equiv}0\;(\mathrm{mod}\;13)$,that
\[
\handwritten{c(7\mathbin{.}13^{n})\equiv0\;(\mathrm{mod}\;13)\quad \text{if}\quad n\geq1}
\tag*{\handwritten{(3\unsure{.|,}3)}}
\]
\[
\handwritten{\text{and}\qquad c(91n)\equiv0\;(\mathrm{mod}\;13)\quad \text{if}\quad(n,7)=1.}
\tag*{\handwritten{(3,4)}}
\]
The first of these results shows that \uline{$c(n)$} is divisible by $13$ infinitely often; the second shows more,that \uline{$c(n)$} is divisible by $13$ in positive density of \uline{$n$}.

% archive sheet number: 43
